\section{A contour shift at the distant cutoff}
\label{sec:perron}

A separate shortening is available where the arithmetic argument estimates
the remainder at its distant cutoff. Because the actual smoothing vanishes
at the endpoint, its Mellin kernel decays quadratically on the initial
vertical line. This permits direct contour inversion. It also gives a
quantitative improvement over the inverse logarithmic powers stated in
\texttt{eventually\_abs\_primeRemainder\_le}.

In this section write $q$ for either modulus $1$ or $Q$ in \eqref{eq:R}, and $d$ for its fixed admissible odd prime. In particular $(d,q)=1$. Write $a_q=1/\varphi(q)$,
$\lambda(n)=\mathbf1_{n\equiv1\bmod q}\Lambda(n)(1-(n/d))$, and
$\mathcal A(s)=\sum_n\lambda(n)n^{-s}$ for $\Re s>1$.
The repository's logarithmic derivative identity gives
$\mathcal A(s)=a_qs/(s-1)+G(s)$. Put $F(s)=G(s)/s$.
The zero-free-region lemma
\texttt{eventually\_norm\_logDerivMean\_div\_le}
in \texttt{New/Reduction/Zeros/PrimeNumberTheorem.lean}
gives, for all sufficiently large $\ell$ depending on $q,d$,
\[
 F \text{ holomorphic},\qquad |F(s)|\le\ell^{26}
 \qquad (\Re s\ge1-\ell^{-26},\ |s|\le e^\ell).
\]
We claim that the centered remainder satisfies
\begin{equation}
 |R_q(x)|\ll_{q,d}\exp\bigl(-\tfrac13(\log x)^{1/27}\bigr).
 \label{eq:perron-improvement}
\end{equation}

Set $U=\log x$ and $\sigma_0=1+1/U$. On this line absolute convergence
and elementary Mellin inversion give
\[
 \sum_{n\le x}\lambda(n)(1/n-1/x)
 =\frac1{2\pi i}\int_{\sigma_0-i\infty}^{\sigma_0+i\infty}
           \frac{\mathcal A(s)x^{s-1}}{s(s-1)}\,ds .
\]
In fact the inverse transform of $y^s/[s(s-1)]$ on a line right of $1$
is $y-1$ for $y\ge1$ and zero for $0<y<1$, by closing the contour
left or right. The $y=1$ value follows by continuity and absolute
convergence. Apply this identity with $y=x/n$ and divide by $x$.
The coefficients obey $0\le\lambda(n)\le2\Lambda(n)$, so
$|\mathcal A(\sigma_0+it)|\ll U$ uniformly in $t$.
This follows either by partial summation from Chebyshev's bound or from
the simple pole of $-\zeta'/\zeta$ at $1$. Truncation at height $T\ge2$
therefore costs $O(U/T)$.

Choose $\ell=U^{1/27}$, $T=e^\ell/2$, and $w=\tfrac12\ell^{-26}$.
Shift the truncated integral to $\Re s=1-w$. The entire closed
rectangle lies strictly inside the stated region for sufficiently large
$x$. Its only pole has residue $a_q U+F(1)$, since the continued
integrand is
\[
        x^{s-1}\left(\frac{a_q}{(s-1)^2}
                       +\frac{F(s)}{s-1}\right).
\]
This is exactly the centering in $R_q(x)$. On the horizontal sides the
length is at most $1$, $|x^{s-1}|\le e$ and $|s-1|\ge T$, giving
$O(\ell^{26}/T)$ for their total. On the left side the absolute integral is
\[
 \ll e^{-wU}\left(w^{-1}+\ell^{26}\log(2+2T/w)\right),
\]
by integrating $(w^2+t^2)^{-1}$ and $(w^2+t^2)^{-1/2}$ for
$|t|\le T$. Here $wU=\ell/2$ and $U=\ell^{27}$, so all three errors
are $O(e^{-\ell/3})$. This proves \eqref{eq:perron-improvement}.
It supplies every fixed inverse logarithmic power needed in the
two-cutoff comparison and dispenses with the gamma-kernel localization
in this application. The narrow classical zero-free region still
provides the shift. The square-root gain at the near cutoff still comes
from the theta chain.
